\originalchapter{AI 辅助研究与新的计算范式}\label{ch:research}
\originalsection{生成器与科学检验}
独立研究不等于独自完成所有操作. 可以让 AI 帮助检索线索、生成候选、写实验代码和提出反例, 但每项工作都需要相应核验. 检索出的题名应回到原文, 代码应实际运行, 证明应逐步检查, 新发现应与已有结果比较. AI 给出“看起来合理”的解释, 只是候选生成的一种产物.

一个常见的有效结构是生成器、评估器与候选档案. 生成器提出程序或构造, 评估器给出明确约束与分数, 档案保留成功、失败和多样性, 下一轮依据档案继续探索. FunSearch 在冻结语言模型与程序搜索的组合中报告数学构造和装箱启发式进展\cite{funsearch}. 这里被改进的是候选程序档案, 基础语言模型权重没有随每轮搜索更新.

AlphaEvolve 的白皮书报告语言模型与演化评估在若干算法任务上的改进, 包括特定 $4\times4$ 复矩阵乘法的 48 次标量乘法构造\cite{alphaevolve}. 这是明确底域、大小与运算模型的结果, 不是任意实矩阵大小都已得到新的渐近复杂度. 作者白皮书、机构部署公告与独立复现也要分别对待. 专门科学系统如 AlphaFold 在蛋白质结构任务的盲评中取得进展\cite{alphafold}, 不意味着蛋白质动力学、功能或药物有效性都已解决.

若问题没有可靠评估器, 整个搜索很容易优化错误代理. 例如只奖励论文摘要读起来新颖, 可能得到更加新颖的虚构. 若反复查看同一个开发集, 候选档案也会对它过拟合. 研究成功必须在冻结选择后的独立任务上验证, 并检查约束、预算和统计误差.

\originalsection{自我改进的三个对象}
第一种是答案级改写: 模型提出答案、检查、重写, 权重不变. 没有外部反馈时, 改写可能只把正确答案改错; 因此应比较首次与改写后的结果, 而不是默认多轮更好. 第二种是系统级改进: 修改代理的提示、代码、工具和搜索策略. Darwin Gödel Machine 维护代理档案并依据编码评估选择代码变体\cite{dgm}. 作者报告的基准提升属于这种系统行为, 不应写成基础模型自行重写所有权重或每步收益已有形式证明.

第三种是参数学习: 生成合成数据、过滤并重新训练模型. 此时必须记录教师来源、数据分布、过滤器、训练成本和测试污染. 如果过滤器只偏爱当前模型已经擅长的答案, 新数据可能减小覆盖面. 如果合成回答含系统性错误, 训练也可能固化它们. 有效反馈与多样化真实证据是关键条件, “自己生成”不是正确性的来源.

所谓递归自我改进, 至少还需证明系统对自身研究能力的改进能够稳定地跨多轮累积, 而不是只提高某个固定评测分数. 现有具体系统为局部可度量改进提供证据; 是否能无限、通用、自动地提升能力仍是远强于这些证据的猜想. 读者应把改进对象和检验范围写出来, 再使用术语.

\originalsection{状态空间与稀疏计算}
稠密注意力保留可直接读取的历史, 代价随序列长度增长. 状态空间模型用有限状态递推, 最简单形式为
\[
h_t=Ah_{t-1}+Bx_t,\qquad y_t=Ch_t.
\]
展开递推得 $h_t=A^th_0+\sum_{j=1}^tA^{t-j}Bx_j$. 因此历史以受矩阵幂控制的滤波形式压入状态. 固定状态维度可使顺序处理随长度线性增长, 但压缩也可能丢掉任意位置精确检索所需的信息. Mamba 引入输入依赖的选择机制, 让状态更新能按内容保留或丢弃信息\cite{mamba}; 不是把这个固定 $A$ 的最简单递推直接当成完整架构.

混合专家模型让路由器只激活部分前馈模块. 总参数数可很大, 每个 token 的激活参数较少. 但路由均衡、容量溢出、通信和推理批次都会影响实际成本. 应同时报告总参数、激活参数与硬件延迟, 不用一个参数数字比较所有架构. 这类稀疏计算改变资源分配, 并不天然保证比相同预算的稠密模型更强.

检索增强把外部文档提供给生成器, 可以让知识更新不必全部经权重训练. 原始 RAG 研究把参数模型和检索的非参数记忆结合\cite{rag}. 检索命中、文档可信、最终回答忠实是三个不同环节. 它既能提供证据, 也增加间接注入和错误来源的入口; 引用链接本身不能保证引用内容支持回答.

\originalsection{扩散与潜空间循环}
自回归不是唯一的生成分解. 离散扩散语言模型先把部分 token 遮蔽或腐化, 再学习根据剩余上下文恢复. 一个基础目标可以写成
\[
\E_{x\sim P,\,m\sim\mathcal Q}\left[-\sum_{i:m_i=1}\log p_\theta(x_i\mid\widetilde x,m)\right],
\]
其中 $m$ 是随机遮蔽集合, $\widetilde x$ 是腐化文本, $\mathcal Q$ 指定不同噪声水平的采样规则. 多轮去噪可以同时改写若干位置, 不同扩散模型有不同目标权重和采样调度. LLaDA 论文探索这类大规模语言生成\cite{llada}; 该简式仅说明去噪机制, 不冒充完整算法. 其质量、步数、可控性和总成本需要与自回归系统按任务比较.

潜空间循环把额外计算放在隐藏状态, 例如 $h_{k+1}=F_\theta(h_k,x)$, 使用共享模块重复更新后再读出答案. 与生成更多思维链 token 相比, 它不要求每次更新都变成文字. 循环越多是否更好取决于动力学与训练. 考虑固定输入和一个非空闭区域 $S$, 假设 $F$ 将 $S$ 映入自身, 且对任意 $u,v\in S$ 满足 $\norm{F(u)-F(v)}\le\rho\norm{u-v}$, $0\le\rho<1$. 从 $h_0\in S$ 开始, 相邻差至多为 $\rho^k\norm{h_1-h_0}$, 所以
\[
\norm{h_{k+m}-h_k}\le\frac{\rho^k}{1-\rho}\norm{h_1-h_0}.
\]
上界随 $k$ 趋于零, 说明任意 $\epsilon>0$ 都可找到 $K$, 使 $k,l\ge K$ 时 $\norm{h_k-h_l}<\epsilon$; 这就是 Cauchy 序列的定义. 本书调用欧氏空间的完备性: 每个这样的序列都收敛. 区域闭保证极限仍在 $S$ 内. Lipschitz 不等式还保证 $F$ 连续, 因而极限 $h_*$ 满足 $F(h_*)=h_*$, 两固定点间距离满足 $d\le\rho d$, 故唯一. 这证明的是状态稳定, 不证明读出答案正确. 一个恒定返回错误答案的递推也完全稳定.

2025 年循环深度研究与 2026 年 STARS 工作研究潜状态推理和更多循环导致的退化, 后者用稳定性约束缓解若干任务上的问题\cite{latent,stars}. 该证据支持进一步比较循环计算路线, 不足以把它宣布为 Transformer 的通用替代或意识机制.

\originalsection{冻结的第一项微研究}\label{sec:micro-study}
第 8 章的一个种子与两个测试组合只能提供起点. 现在实际检验一个更明确的问题: 在完整文档组合留出时, 小模型的叫声预测是否只靠更近的动物, 还是符合句首主语依赖? 本节保留同一架构、步数与优化器, 改成五个初始化和三个预先冻结的划分. 结果出炉前确定问题、指标、基线与干预, 避免看完测试后换假说、挑种子.

数据仍只有八个完整文档, 一律形如 \texttt{<bos> speaker sees color other says sound <eos>}. 主语与干扰动物各取 cat/dog, 颜色取 red/blue, 主语决定 meow/woof. 按“主语, 颜色, 干扰动物”的循环顺序编号 $0,\ldots,7$; 例如文档 0 为 cat/red/cat, 文档 3 为 cat/blue/dog, 文档 4 为 dog/red/cat. 叫声是规则给定的标签, 不是从网络观察中猜出的因果规律.

三个固定划分如下, 每个取余下六文档训练. A 与 B 的测试主语和干扰动物不同, C 相同. A/B 的训练中有四个匹配、两个失配, C 则两个匹配、四个失配. 测试不与训练重复, 但三个划分互相共享文档; 不能把它们视为三个独立随机测试集.
\begin{center}\begin{tabular}{lll}\toprule
划分 & 留出编号 & 测试主语/颜色/干扰动物\\\midrule
A & 3,4 & cat/blue/dog; dog/red/cat\\
B & 1,6 & cat/red/dog; dog/blue/cat\\
C & 0,7 & cat/red/cat; dog/blue/dog\\\bottomrule
\end{tabular}\end{center}

种子固定为 144--148, 每个划分各训练五次. 每次新建模型、Adam 一阶与二阶状态清零; 相同种子在不同划分使用相同初始权重, 参数哈希在训练前记录. 共 15 次全批次更新实验, 均为单头 $D=16$, 前馈宽度 32, 2,538 参数, 1,400 步, 学习率 0.01, 梯度总范数裁剪为 1. 没有架构搜索、测试调参或反复重跑直到成功. 冻结 JSON 与运行前澄清的哈希记录在源码 \texttt{research/} 下; 所有结果而非最好种子写入 JSON.

两个简单基线只预测 \texttt{says} 后的叫声. 最近动物规则把 \texttt{other} 视作主语, A/B 原输入准确率为 0, C 为 1. 最后词频规则估计训练中 \texttt{says} 后各叫声频率; 每个训练集各有三次 meow/woof, 概率都是 $1/2$, 平票固定输出 meow, 三个测试准确率都为 $1/2$. 这些基线没有完整序列模型, 所以不能把它们与网络的全 token NLL 横向比较.

\originalsection{配对干预与误差}
保持其他内容不变, 第一项干预翻转干扰动物并保留叫声标签; 第二项翻转句首主语并据规则改写叫声标签. 评价在输入位置 5 的 \texttt{says} 处进行, 因果掩码只允许它读位置 0--5, 不能读后面的正确叫声. 完整新文档中的叫声也同步修改, 但本节不计算干预文档的全序列 NLL. 两项干预在所有六个原测试文档上都回到对应训练组合, 所以它们测试局部输入敏感性, 不是另加的未见组合泛化测试.

预定预测是: 若叫声稳定依赖主语, 改干扰动物后应保持, 改主语后应翻转. 只看干预准确率会误导: 最近动物基线在 A/B 的两类干预中也达到 100\%, 因为翻转任一动物都把原先失配变成匹配或把标签改成干扰动物的叫声; 在 C 两类干预反而均为 0. 必须把原输入、两项干预和基线一起比较.

对每个模型与原测试文档, 同时记录预测、正确叫声概率、原叫声概率和训练成员状态. 用配对差
\[
\Delta_{\rm speaker}=p_{\rm speaker\ flip}(y_{\rm original})-p_{\rm original}(y_{\rm original}),\qquad
\Delta_{\rm other}=p_{\rm other\ flip}(y_{\rm original})-p_{\rm original}(y_{\rm original})
\]
描述变化. 主语翻转有了新正确答案, 这里仍跟踪旧答案概率, 所以合理预测是第一项明显为负, 第二项接近零. 它们不能合并成一个“干预分数”.

全部 15 次训练, 原输入与两类干预的叫声准确率都为 1. 下表训练/测试 NLL 为各划分五个初始化的描述性平均, 全对次数则是一次初始化的模型是否答对该划分的两个叫声. 区间的严格条件在下节说明.
\begin{center}\begin{tabular}{lrrrrr}\toprule
划分 & 全对次数 & 叫声准确率 & 训练 NLL & 测试 NLL & 95\% 区间\\\midrule
A & 5/5 & 1.000 & 0.255990 & 1.474148 & $[0.4782,1]$ \\
B & 5/5 & 1.000 & 0.256043 & 1.450169 & $[0.4782,1]$ \\
C & 5/5 & 1.000 & 0.255998 & 1.440073 & $[0.4782,1]$ \\
\bottomrule\end{tabular}\end{center}

\begin{center}\begin{tabular}{lrrr}\toprule
种子 & A 测试 NLL & B 测试 NLL & C 测试 NLL\\\midrule
144 & 1.506582 & 1.476743 & 1.473264 \\
145 & 1.538308 & 1.161478 & 1.434869 \\
146 & 1.360304 & 1.496133 & 1.506396 \\
147 & 1.543225 & 1.531221 & 1.288358 \\
148 & 1.422322 & 1.585268 & 1.497475 \\
\bottomrule\end{tabular}\end{center}

30 个配对文档上, 主语替换的 $\Delta_{\rm speaker}$ 落在 $[-0.999993,-0.992537]$, 干扰替换的 $|\Delta_{\rm other}|$ 最大为 $0.007446$. 这些是全部预定案例的描述量, 不是总体效应区间.


这组结果与主语依赖预测相符, 原输入表现排除了单纯最近动物或最后词频解释. 但它没有唯一识别内部算法: 模型可能结合位置、词表或多条路径完成这个小函数. 干预回到训练组合, 不控制自然语言中的共指、句法或世界知识. 不能由此声称发现了通用推理、某个注意力头的因果电路或独立研究能力.

完整测试 NLL 在 1.44--1.47, 训练约为 0.256; 叫声全对没有消除其他位置的失配. 把每个位置的负对数概率分别平均, 再对位置取平均就恢复完整 NLL. 五个种子的逐位置平均如下, 每列是所预测的目标 token:
\begin{center}\begin{tabular}{lrrrrrrr}\toprule
划分 & 主语 & sees & 颜色 & 干扰动物 & says & 叫声 & EOS\\\midrule
A & 0.6932 & 0.0000 & 1.0986 & 8.5272 & 0.0000 & 0.0001 & 0.0000 \\
B & 0.6933 & 0.0000 & 1.0991 & 8.3587 & 0.0000 & 0.0001 & 0.0000 \\
C & 0.6932 & 0.0000 & 1.0979 & 8.2886 & 0.0000 & 0.0008 & 0.0000 \\
\bottomrule\end{tabular}\end{center}


主要损失来自干扰动物, 不是叫声. 被留出的组合使训练中的某些“主语+颜色”前缀后只出现一种干扰动物. 如 A 中 cat/blue 前缀后只训练过 cat, 测试却要求 dog. 交叉熵拟合训练条件分布时会把未出现的分支概率压得很小, 测试负对数概率因此很大. 六文档还使每个主语下的测试颜色在训练中仅占 $1/3$, 所以颜色 NLL 接近 $-\log(1/3)=\log3$, argmax 偏向另一颜色. 这部分不能简单叫作确定规则没有学会: 颜色本来是随机分支, 而留出规则改变了测试中的分支比例.

下一个研究应把“叫声规则”和“完整生成分布”分开评价, 另冻长度与组合生成器, 再检验未见前缀和分支校准. 当前错误分析没有据此挑新架构或学习率, 也不把观察到的误差用于重写原先成功指标. 可以保留假说得到局部支持的结果, 同时完整报告其未达到的目标.

\originalsection{样本、划分与初始化的不确定性}\label{sec:micro-ci}
“重复五次”必须说明重复的对象. 若抽取新测试文档, 随机性来自文档分布; 若重新选训练/测试划分, 来自数据使用方式; 若数据固定但重新初始化, 来自训练过程. 本研究只有第三种被重复. 文档是八个固定组合中的指定留出, 三个划分也是事先指定并互相重叠, 没有对任意新文档或任意划分进行概率抽样.

固定一个划分, 定义 $I_s=1$ 当第 $s$ 次初始化所得模型答对两个测试叫声, 否则为 0. 假设训练配置不变、初始化独立同分布, 每次该事件概率为 $p$, 则 $K=\sum_{s=1}^n I_s$ 满足
\[
\Pr_p(K=k)=\binom nk p^k(1-p)^{n-k},\qquad k=0,\ldots,n.
\]
理由是任意指定的 $k$ 次成功事件有上述乘积概率, 共有 $\binom nk$ 种不同位置集合, 它们互斥. $p$ 指在这个固定划分上的\emph{整次训练全对概率}, 不是随机一个 token 的准确率, 更不是自然语言研究能力概率. 五个种子只是公开且可复现的初始化选择; 统计覆盖率解释依赖所假设的初始化抽样方式, 不能把固定编号自动当成实际随机采集的人口样本.

观测比例记为 $\hat p=K/n$. 少样本全对时, 常见插值标准误 $\sqrt{\hat p(1-\hat p)/n}$ 为零, 会错误给出 $[1,1]$. 用二项分布尾概率反演可避免这个边界问题\cite{nist-ci}. 设显著性水平 $\alpha=0.05$. 对 $0<k<n$, 下界 $L$ 与上界 $U$ 分别解
\begin{align*}
\sum_{j=k}^{n}\binom njL^j(1-L)^{n-j}&=\alpha/2,\\
\sum_{j=0}^{k}\binom njU^j(1-U)^{n-j}&=\alpha/2.
\end{align*}
第一式问: 若真实 $p$ 更小, 观察到至少 $k$ 次成功是否已太罕见? 第二式问: 若 $p$ 更大, 至多 $k$ 次是否太罕见? 为说明求根确实可行, 记两个左端尾多项式为 $F_k(p)$ 与 $H_k(p)$. 对求和式逐项求导, 用 $j\binom nj=n\binom{n-1}{j-1}$ 与 $(n-j)\binom nj=n\binom{n-1}{j}$ 移动指标, 两串和中间项相消, 得
\[
F_k'(p)=n\binom{n-1}{k-1}p^{k-1}(1-p)^{n-k}>0,\qquad
H_k'(p)=-n\binom{n-1}{k}p^k(1-p)^{n-k-1}<0.
\]
这里 $0<k<n$, $0<p<1$. 此外 $F_k(0)=0,F_k(1)=1$, $H_k(0)=1,H_k(1)=0$. 连续性与严格单调保证对 $\alpha/2\in(0,1)$ 各有唯一根, 程序用二分法. 端点 $k=0$ 取 $L=0$, $k=n$ 取 $U=1$, 其另一端仍用相应尾式求出.

覆盖保证来自拒绝两端过于罕见的计数. 对任意固定真值 $p$, 事件 $L(K)>p$ 等价于观察值落在上尾概率小于 $\alpha/2$ 的计数区; 这些计数是一个连续上尾, 总概率至多 $\alpha/2$. 同理 $U(K)<p$ 的概率至多 $\alpha/2$. 两事件相加给 $\Pr_p\{p\notin[L(K),U(K)]\}\le\alpha$. 因为计数离散, 实际覆盖常超过 $1-\alpha$. 这是抽样重复时区间覆盖固定真值的性质, 不是观察数据后“$p$ 在本区间中的后验概率为 95\%”.

本次每个划分 $n=k=5$, 下界满足 $L^5=0.025$, 所以 $L=0.025^{1/5}\approx0.478176$, 上界为 1. 5/5 的整齐结果仍只得到宽区间 $[0.4782,1]$. 三个区间分别具有上述条件性的点态解释, 并未作三者同时覆盖校正. 它们也不适用于平均准确率、NLL 或概率差. 同种子跨划分的权重相同, 结果配对; 若描述跨三个固定划分的平均, 应先在每个种子内平均, 不能当作 15 次独立同分布试验再缩小误差条.

两个固定测试文档及其干预、同文档 token 均有结构依赖. 将 30 个叫声预测重新抽样做 IID bootstrap 不会创造 30 份独立数据, 本项目未这样做. 若后续研究改用随机文档, 要先说明独立单位是文档、任务还是生成模板, 再按这个层级设计估计. 当前不估计任意新划分的不确定性, 这个缺口由新增实验设计解决, 无法用已有数据的重抽样补齐.

实验选择、种子与主要参数有公开命令入口. 先复现参考结果, 另选输出目录避免覆盖, 再只改一个因素并预先写下预测:
\begin{verbatim}
python experiments/run_all.py --experiment attention_walkthrough \
  --out /tmp/attention-check
python experiments/run_all.py --experiment residual_debug \
  --seed 19 --out /tmp/residual-check
python experiments/run_all.py --experiment micro_research \
  --out /tmp/frozen-study
python experiments/run_all.py --experiment language_model \
  --seed 145 --steps 200 --lr .01 --width 8 \
  --out /tmp/exploratory-lm
\end{verbatim}
最后一条是探索配置, 不冒充本节冻结研究. 微研究也允许 \texttt{--seeds}、\texttt{--steps}、\texttt{--lr}、\texttt{--width} 与 \texttt{--design} 指定另外方案, 参数覆盖会写入 \texttt{exploratory\_overrides}. 原始概率、训练曲线、初始化与源码哈希、CLI 参数和环境留在结果目录. 耗时属于机器测量, 不用较短耗时外推大模型成本.

另一可扩展项目是八动作的偏好学习. 先由已知奖励得到解析 KL 最优策略, 再生成带噪偏好训练 DPO, 比较真实奖励、KL 与目标差距. 这让后训练不再只是术语, 因为真实效用与偏好拟合可以分开测量. 本项目已经提供并实际运行这项有限动作实验; 它不训练大语言模型.

一个月的可执行研究节奏可以按产物安排: 第一阶段复算数学、跑通固定基线; 第二阶段冻结生成器、划分和指标; 第三阶段执行预定干预、保留全部原始结果; 第四阶段写出结论、反例、限制和复现实验命令. 如果所有机制检验失败, 一份说明失败位置的结果仍是有效研究笔记, 不必为获得“发现”而改写事实.

\originalsection{练习与解答}
\begin{exercise}
如果代理不断修改自身提示, 但基础模型参数不变, 应归在哪一类改进? 要证明有效还缺什么?
\end{exercise}
\begin{solution}
它属于系统级改进, 也可能包含答案级改写. 要在冻结选择后的独立测试上与相同预算的固定提示和随机变体搜索比较, 保留训练/开发反馈和失败记录. 单看反复用过的开发集分数上升不足以证明泛化改进.
\end{solution}
\begin{exercise}
递推 $h_{k+1}=0.5h_k+1$ 从任意实数初值收敛到什么? 若读出器将此状态解释为错误答案, 收敛还能证明推理正确吗?
\end{exercise}
\begin{solution}
固定点方程为 $h=0.5h+1$, 得 $h_*=2$. 误差满足 $h_k-2=0.5^k(h_0-2)$, 故收敛. 读出语义不由收缩性控制, 因而不能据此证明推理正确. 稳定性与任务有效性应分别评价.
\end{solution}
\begin{exercise}
在一个程序搜索任务中, 为什么应禁止候选程序修改评分器和隐藏测试? 允许修改会改变哪个研究对象?
\end{exercise}
\begin{solution}
若允许候选修改评分器, 高分可能来自降低标准或改写测试答案, 而不是算法改善. 此时研究对象已从固定规格下的候选性能变为整个评价系统的变动, 原来的分数不再可比较. 可信评分器和测试数据应置于候选不可修改的边界之外.
\end{solution}
\begin{exercise}
五个初始化在固定两个文档上全对. 有人报告“十个独立样本的准确率为 100\%, 标准误为零”; 另有人把三个划分合成 15 次全对, 使用 $0.025^{1/15}$ 为下界. 指出两处单位错误, 并给出本研究可解释的区间.
\end{exercise}
\begin{solution}
同一模型的两个文档不是独立初始化, 标签分支也由指定组合选择. 本节事件是“一次训练两个都对”, 所以每个划分只有五次. 三个划分复用初始化且文档重叠, 未满足 15 次 IID Bernoulli 假设. 对每个固定划分分别报告条件性的 $[0.025^{1/5},1]\approx[0.4782,1]$, 不声称同时覆盖或自然语言泛化. 全对时经验方差零不表示真实概率已确定为 1.
\end{solution}
\label{ch:end:13}
